\chapter{Per-Transfer Decisions, Compression and Storage}
\label{ch:deciding}

The pool of the previous chapter offers two paths, a quantizer offers a smaller
encoding, and a storage tier offers a location off the GPU, so a choice among
them exists for every transfer. Five components make that choice or carry
it out. The first three decide: a
classifier that labels a transfer hot or cold, a rule tree that turns labels
and system state into a decision, and a controller that was intended to
balance the two queues and has an issue in its switching rule, which we
analyze. The last two act
on a decision: the quantizer packs cache vectors into fewer bytes for the
wire and for disk, and the storage tier moves cache blocks to an NVMe drive
when GPU memory is short. Neither changes the cache with which the model
computes, which remains in its 16-bit format on the GPU. Of these components, the static classifier and the quantizer run on the
serving path of \cref{fig:architecture}, and the others are evaluated on
their own, so each section states what was measured and on what data. Each
component is presented in the same order: what it does, the algorithm, and
the properties that we can prove about it.

\section{Observations Available to the Middleware}
\label{sec:observations}

A decision needs inputs, and the middleware collects them at the one place
through which every transfer passes, the sender. Every transfer produces a
small record: the logarithm of its size, how often
this tensor has been accessed, the queue depth when it was submitted, whether
the request is in prefill or decode, the logarithm of the sequence length, the
logarithm of the time since the previous transfer, and whether it is a control
message. Latencies go into histograms per class, and events such as timeouts,
reconnections and fallbacks to TCP are counted. Everything is exported in
Prometheus text format, and the records can be written to a log that the
\looseness=-1 learned classifier trains on.

\section{A Classifier as a Weighted Automaton}
\label{sec:classifier}

A weighted finite automaton is a finite set of states with transitions that
carry weights instead of being simply allowed or forbidden
\citep{droste2009wfa,suresh2021wfa}. The classifier keeps one such automaton
per tensor, with three states (\cref{fig:wfa}). A tensor in the hot state is
sent on the hot endpoint.

\begin{figure}[htbp]
  \centering
  \begin{tikzpicture}[diagram,
      alt={State diagram of the classifier with three states in a row: cold, warm and hot.
        Arrows lead from cold to warm and from warm to hot, labeled used often or small.
        Arrows lead back from hot to warm and from warm to cold, labeled idle or large. Each
        state also has a loop back to itself. A note says every arrow carries a weight computed
        from the features of the transfer, and that transfers in the hot state use the hot
        endpoint.},
      st/.style={tint=#1, circle, minimum size=17mm, font=\footnotesize}]
    \node[st=dBlue]   (c) at (0,0)    {Cold};
    \node[st=dPurple] (w) at (5.6,0)  {Warm};
    \node[st=dOrange] (h) at (11.2,0) {Hot};
    \draw[flow] (c) to[bend left=22] node[midway, above=2pt, small text] {used often, or small} (w);
    \draw[flow] (w) to[bend left=22] node[midway, above=2pt, small text] {used often, or small} (h);
    \draw[flow] (h) to[bend left=22] node[midway, below=2pt, small text] {idle, or large} (w);
    \draw[flow] (w) to[bend left=22] node[midway, below=2pt, small text] {idle, or large} (c);
    \draw[flow] (c) to[loop left, looseness=5] (c);
    \draw[flow] (w) to[loop below, looseness=5] (w);
    \draw[flow] (h) to[loop right, looseness=5] (h);
    \node[small text, align=center] at (5.6,-2.75) {Every arrow has a weight computed from the transfer's features.\\Transfers in the hot state use the hot endpoint.};
  \end{tikzpicture}
  \caption[The three-state automaton kept for each tensor]{The three-state automaton kept for each tensor.}
  \label{fig:wfa}
\end{figure}

\subsection{The Static Version and Its Limitation}

The static classifier uses fixed thresholds (\cref{alg:static}). A transfer
under 4 KiB is hot, one over 256 KiB is cold, and one in between becomes hot
once its tensor has been accessed ten times. A phase detector then adjusts
the class: it examines the layer numbers of the last eight transfers, and it
reports prefill if they increase by one each time, decode if the last three
are equal, and unknown otherwise. In prefill a hot transfer over 2 KiB is
moved to the cold class, and in decode a cold transfer is promoted once its
tensor has been accessed five times.

However, the decode rule of the detector never fires on the traffic for which
it is intended.

\begin{proposition}[The static detector never reports decode]
\label{prop:nodecode}%
Let $L \ge 2$ be the number of layers of the model, and suppose that
transfers arrive layer by layer, with the layer numbers $0, 1, \dots, L-1$
repeating. Then the phase detector never reports decode.
\end{proposition}

\begin{proof}
The detector reports decode only when the last three layer numbers are equal.
In the stated traffic, consecutive transfers always have different layer
numbers: the next number is either one higher, or it is $0$ after $L-1$, and
$0 \ne L-1$ because $L \ge 2$. Hence no two consecutive layer numbers are equal,
and no three are.
\end{proof}

\begin{pseudocode}{Static classification of one transfer}{alg:static}
  \Function{Classify}{$\var{key}$, $\var{size}$, layer $\ell$}
    \State $n \gets$ number of transfers seen under $\var{key}$, this one included
    \State $\var{phase} \gets \Call{Phase}{\ell}$
    \If{$\var{size} < 4$ KiB}
      \State $c \gets$ hot
    \ElsIf{$\var{size} > 256$ KiB}
      \State $c \gets$ cold
    \Else
      \State $c \gets$ hot if $n \ge 10$, otherwise cold
    \EndIf
    \If{$\var{phase}$ is prefill, $c$ is hot and $\var{size} > 2$ KiB}
      \State $c \gets$ cold
    \ElsIf{$\var{phase}$ is decode, $c$ is cold and $n \ge 5$}
      \State $c \gets$ hot
    \EndIf
    \State \Return $c$
  \EndFunction
  \Function{Phase}{$\ell$}
    \State append $\ell$ to a window that keeps the last eight layer numbers
    \State \textbf{if} the window holds fewer than three numbers \textbf{then return} the phase reported last
    \State \textbf{if} every number is one more than the number before it \textbf{then return} prefill
    \State \textbf{if} the last three numbers are equal \textbf{then return} decode
    \State \Return unknown
  \EndFunction
\end{pseudocode}

Decode traffic is exactly of this form: each generated token touches every
layer in order. A window of eight that lies inside one pass over the layers
has differences of one and is labeled prefill. A window that contains the
wrap from $L-1$ to $0$ is labeled unknown. With 40 layers, 7 windows in
every 40 contain the wrap, so a decode-only stream should be labeled prefill
82.5 percent of the time, unknown 17.5 percent, and decode never. A synthetic
stream of 8,000 decode transfers gave 6,605 prefill, 1,395 unknown and 0
decode, and the one recorded run on a real device gave no decode labels
either (\cref{sec:res-loopback}).

The consequence is that the cache updates of a running request, which are
small, urgent, and in the 4 KiB to 256 KiB band, are treated as bulk until
each has been seen ten times. On the serving path of \cref{sec:overview} the
effect is more direct, as every blob carries a new key and its count never
rises above one.

\subsection{The Learned Version}

The replacement keeps the three states and learns the weights. For a transfer
$x$ whose tensor is in state $s$, the probability of moving to state $s'$ is
a softmax over the three possible next states,
\mathalt{the probability of s prime given s and x equals the exponential of W s s prime dot
  phi of x, divided by the sum over the three states r of the exponential of W s r dot phi
  of x}{%
\begin{equation}
  P(s' \mid s, x) = \frac{\exp\!\big(W_{s,s'} \cdot \phi(x)\big)}{\sum_{r} \exp\!\big(W_{s,r} \cdot \phi(x)\big)},
  \label{eq:wfa}
\end{equation}}
where $\phi(x)$ is the vector of features of the transfer, $W_{s,s'}$ is a
vector of weights of the same length, one for each pair of a current state
$s$ and a next state $s'$, the dot is the dot product of \cref{eq:dot}, and
the sum in the denominator runs over the three states $r$. The features include
the two phase indicators and the products of each with the logarithm of the
size, so the model can hold a different size cut-off in each phase, and a
different sharpness around it. Training minimizes cross-entropy against
labeled transitions with a quadratic penalty on the weights. Two properties
are proved in \cref{sec:proofs-wfa}: the training objective is strictly convex
and has a unique minimum, and a per-phase size cut-off of the kind the static
classifier lacks can be represented exactly.

\subsection{Evaluation on Synthetic Traffic}

The classifier was evaluated on generated traffic with a rule-based oracle. \Cref{tab:wfa} gives the fraction
of transfers routed as the oracle routes them, on a model shape that was not
\looseness=-1 in the training set.

\begin{table}[htbp]
  \centering
  \caption[Routing accuracy on synthetic traffic, 4,000 held-out events per mix]{Routing accuracy on synthetic traffic, 4,000 held-out events per mix. Synthetic data with a synthetic oracle; not a measurement of real traffic.}
  \label{tab:wfa}
  \small
  \begin{tabular}{L{4.6cm}L{3.2cm}L{3.2cm}}
    \toprule
    Traffic mix & Static thresholds & Learned weights \\
    \midrule
    Balanced & 0.261 & 0.999 \\
    Decode-heavy & 0.171 & 0.999 \\
    Prefill-heavy & 1.000 & 0.999 \\
    Bursty & 0.489 & 0.999 \\
    \bottomrule
  \end{tabular}
\end{table}

The static classifier is correct when the traffic is prefill and incorrect in proportion to the amount of decode traffic, as \cref{prop:nodecode} predicts.
The learned classifier's score is partly by construction, because the oracle
that labels the test data is the same kind of rule that generated the
training data. What the table shows is that the model can represent the
intended policy and that training finds it. Whether it is the right policy for
real traffic requires traces labeled by measured latency, which is future
work.

\subsection{Analysis of the Learned Version}
\label{sec:proofs-wfa}

The notation is the following. The feature vector $\phi(x)$ of a transfer
$x$ has $F$ components. It consists of the standardized features of
\cref{sec:observations}, among them the logarithm of the size, written $t$,
and the two phase indicators, $p$ for prefill and $e$ for decode, each of
which is $1$ in its phase and $0$ otherwise, together with the products $t p$
and $t e$ and a final component that is the constant $1$. For each current
state $s$ and next state $s'$ there is a weight vector $W_{s,s'}$ with $F$
components, the score of the next state $s'$ is $W_{s,s'} \cdot \phi(x)$, and
$W$ stands for all weight vectors together. The training set has $n$ examples
$(x_i, s_i, y_i)$, each of which is a transfer $x_i$, the previous state
$s_i$ of its tensor, and the labeled next state $y_i$. The previous state is
taken from the previous label of the same tensor and not from the model's own
output. The objective is
\mathalt{J of W equals minus 1 over n times the sum over i of the logarithm of the probability of
  y i given s i and x i, plus eta over 2 times the squared length of W}{%
\begin{equation}
  J(W) = -\frac{1}{n} \sum_{i=1}^{n} \log P(y_i \mid s_i, x_i)
         + \frac{\eta}{2}\, \lVert W \rVert^2 , \qquad \eta > 0,
  \label{eq:wfa-loss}
\end{equation}}
where $P$ is the probability of \cref{eq:wfa}, $\lVert W \rVert^2$ is the sum
of the squares of all weights, and $\eta$ is the weight of the penalty.

\begin{proposition}[The training objective has one minimum]
\label{prop:convex}%
The objective $J$ of \cref{eq:wfa-loss} is strictly convex in $W$ and has
exactly one minimizer.
\end{proposition}

\begin{proof}
\emph{The loss of one example.} Write $z_{i,r} = W_{s_i,r} \cdot \phi(x_i)$
for the score of next state $r$ on example $i$. Each score is a linear
function of $W$, and by \cref{eq:wfa} the loss of example $i$ is
\mathalt{minus the logarithm of the probability of y i equals minus z i y i, plus the logarithm
  of the sum over r of the exponential of z i r}{%
\begin{equation}
  -\log P(y_i \mid s_i, x_i) = -z_{i,y_i} + \log \sum_{r} \exp(z_{i,r}).
\end{equation}}
The first term is linear in $W$, and the second is the log-sum-exp function
of the three scores.

\emph{Log-sum-exp is convex.} Its matrix of second derivatives is
$\operatorname{diag}(\pi) - \pi \pi^{\mathsf T}$, where
$\pi = (\pi_1, \pi_2, \pi_3)$ is the vector of the three probabilities of
\cref{eq:wfa}, $\operatorname{diag}(\pi)$ is the diagonal matrix with these
entries, and $\pi^{\mathsf T}$ is the transpose of $\pi$. For any vector
$v = (v_1, v_2, v_3)$,
\mathalt{v transposed times the quantity diag of pi minus pi pi transposed, times v, equals the
  sum over r of pi r times v r squared, minus the square of the sum over r of pi r times v r,
  which is at least zero}{%
\begin{equation}
  v^{\mathsf T} \big(\operatorname{diag}(\pi) - \pi \pi^{\mathsf T}\big)\, v
  = \sum_r \pi_r v_r^2 - \Big(\sum_r \pi_r v_r\Big)^{2} \ge 0 ,
\end{equation}}
because the right side is the variance of the values $v_r$ under the
probabilities $\pi_r$. The matrix is therefore positive semidefinite, and
log-sum-exp is convex.

\emph{The objective is strictly convex.} A convex function of a linear
function of $W$ is convex, and a sum of convex functions is convex, so the
first term of \cref{eq:wfa-loss} is convex in $W$. The penalty has the
second-derivative matrix $\eta I$, where $I$ is the identity matrix. This
matrix is positive definite, so the penalty is strictly convex, and the sum
\looseness=-1 of a convex and a strictly convex function is strictly convex.

\emph{A minimizer exists and is unique.} A strictly convex function has at
most one minimizer. One exists, because the loss of each example is the
negative logarithm of a probability and so is not negative, hence
$J(W) \ge \tfrac{\eta}{2} \lVert W \rVert^2$, which grows without bound as $W$
does, and $J$ is continuous.
\end{proof}

The consequence is that the weights training converges to do not depend on
where it starts. The implementation runs a fixed number of steps of a
gradient method, so it approaches the minimizer and is not guaranteed to
reach it exactly; the notebook checks the loss curve.

\begin{proposition}[A size cut-off per phase is representable]
\label{prop:cutoff}%
Fix a current state $s$ and two next states $s'$ and $s''$. Let $\tau_P$ and
$\tau_D$ be two cut-offs on the logarithm $t$ of the size, one for prefill
and one for decode, and let $\gamma_P$ and $\gamma_D$ be two positive slopes.
Then there are weights for which the score of $s'$ minus the score of $s''$
is $-\gamma_P (t - \tau_P)$ for every prefill transfer and
$-\gamma_D (t - \tau_D)$ for every decode transfer. With those weights the
model prefers $s'$ to $s''$ exactly when $t < \tau_P$ in prefill and exactly
when $t < \tau_D$ in decode.
\end{proposition}

\begin{proof}
\emph{The target is linear in the features.} The difference of scores to be
produced is
\mathalt{Delta of x equals minus gamma P times the quantity t p minus tau P p, minus gamma D
  times the quantity t e minus tau D e}{%
\begin{equation}
  \Delta(x) = -\gamma_P\, (t p - \tau_P\, p) \;-\; \gamma_D\, (t e - \tau_D\, e),
\end{equation}}
a linear combination of the four raw features $t p$, $t e$, $p$ and $e$.
Standardizing a feature subtracts a constant and divides by a positive
constant, so each raw feature is a linear function of its standardized form
and the constant feature. Hence there is a vector $v$ with
$v \cdot \phi(x) = \Delta(x)$ for every transfer $x$.

\emph{The weights.} Set $W_{s,s'} = v$ and $W_{s,s''} = 0$, so that the score
of $s'$ minus the score of $s''$ is $\Delta(x)$. In prefill $p = 1$ and
$e = 0$, so the difference is $-\gamma_P (t - \tau_P)$, which is positive
exactly when $t < \tau_P$. In decode $p = 0$ and $e = 1$, so the difference
is $-\gamma_D (t - \tau_D)$, which is positive exactly when $t < \tau_D$.
\end{proof}

Two remarks. The phase indicators alone would already allow two different
cut-offs; the products add the freedom to make the two slopes differ, which
sets how sharply the probability changes around each cut-off. Moreover, the proposition states that the intended rule is available to training, not that
training must find it: the minimizer of \cref{prop:convex} depends on the
labels. The static classifier cannot hold a decode cut-off at all, because
its decode rule waits for a detector that never fires
(\cref{prop:nodecode}).

\section{A Rule Tree for the Decision}
\label{sec:policy}

A classifier answers one question, hot or cold, from the history of one
tensor, whereas a decision also depends on state that no single transfer
shows, such as how full the GPU memory is and how busy the link is. The rule
tree combines the two, so that the label of the classifier is one input among
several. The decision for a
transfer has two fields, a route (hot, cold or spill to storage) and whether
to compress. A policy is a list of rules, each with a name, a priority, a
condition and an action that sets one or both fields. A condition is a tree
whose inner nodes are and, or, and not, and whose leaves compare a field of
the transfer or of the system state with a constant, test membership in a
set, or ask the classifier. The form follows the rule composition of a
request router built for the same serving engine
\citep{liu2026router,wang2025router}.

\begin{pseudocode}{Evaluating a policy for one transfer}{alg:policy}
  \State $\var{decision} \gets$ empty
  \For{each rule, from highest priority to lowest}
    \If{the rule's condition holds for this transfer}
      \State set each field the rule names, if that field is still unset
      \If{the rule is marked final, or both fields are set}
        \State \textbf{stop}
      \EndIf
    \EndIf
  \EndFor
  \State set any field still unset to its default: route cold, no compression
  \State \Return $\var{decision}$
\end{pseudocode}

Three properties hold by construction and are proved in \cref{sec:proofs-policy}:
each field is set by the highest-priority rule that matches and names it;
evaluation may stop early without changing the result; and the cost is linear
in the size of the rule set. A field that a condition refers to and the
transfer lacks makes that leaf false and is recorded, so a misspelled field
name shows up in the explanation of a decision instead of silently matching
nothing.

\Cref{tab:policy} is the default policy.

\begin{table}[htbp]
  \centering
  \caption[The default policy]{The default policy. Higher priority is evaluated first.}
  \label{tab:policy}
  \small
  \begin{tabular}{L{1.5cm}L{7.6cm}L{4.4cm}}
    \toprule
    Priority & Condition & Action \\
    \midrule
    100 & control message & route hot \\
    90 & GPU memory at least 90 percent full, size over 256 KiB, not decode & route to storage \\
    80 & decode phase and a small cache update & route hot \\
    70 & prefill phase and size over 64 KiB & route cold \\
    60 & link at least 80 percent busy, route cold or storage, size at least 256 KiB & compress \\
    30 & classifier says hot and size at most 256 KiB & route hot \\
    20 & size under 4 KiB & route hot \\
    \bottomrule
  \end{tabular}
\end{table}

The thresholds of the default policy are initial values taken from the sizes
in \cref{tab:traffic}. Setting them from measured traffic is future work, and no reported result
depends on them.

\subsection{Properties of the Rule Tree}
\label{sec:proofs-policy}

The procedure is \cref{alg:policy}. Rules are sorted once, by priority from
high to low with ties kept in file order, and the decision starts with both
fields unset. The condition of a rule is evaluated on the transfer, the
system state and the decision so far, and if it holds, the rule sets each
field that its action names and that is still unset. Evaluation stops after a
rule marked final or when both fields are set, and fields that are still
unset then take their defaults.

\begin{proposition}[Each field belongs to one rule]
\label{prop:owner}%
When the procedure ends, each field holds the value given by the first rule,
in sorted order, that was evaluated, whose condition held, and whose action
names that field. If there is no such rule, the field holds its default.
\end{proposition}

\begin{proof}
The proof is by induction on the number of rules processed.

\emph{Invariant.} Every field that is set holds the value from the first
processed rule that matched and names it, and for every field that is unset
no processed rule matched and names it.

\emph{Before the first rule.} Nothing is set, and the invariant holds
trivially.

\emph{One more rule.} Processing a rule changes nothing if its condition
fails. If the condition holds, the rule sets only fields that are unset. By
the invariant no earlier rule matched and names those, so this rule is the
first for them, and fields already set are not touched.

\emph{At the end.} Stopping keeps the invariant, and the defaults fill only
fields that no processed rule set.
\end{proof}

\begin{proposition}[Stopping early does not change the decision]
\label{prop:earlystop}%
The decision is the same as if every node of every rule were evaluated, in
two senses. Inside a condition, an and node may stop at its first false
child and an or node at its first true child. Across rules, evaluation may
stop once both fields are set.
\end{proposition}

\begin{proof}
\emph{Inside a condition.} For an and node, one false child makes the node
false whatever the others are, and if no child is false, all were evaluated.
The same holds for an or node with true. This is safe only if evaluating a
child cannot change what another child returns. Comparison and membership
leaves read the transfer, the system state and the decision so far, none of
which changes during one condition. A classifier leaf could be different,
because a classifier may keep history. The engine therefore calls each named
classifier at most once per decision and reuses the answer at every leaf
that names it, so every leaf returns the same value whenever and however
often it is evaluated.

\emph{Across rules.} A rule can set only fields that are unset. Once both
fields are set no later rule can change the decision, so not evaluating the
later rules changes nothing.
\end{proof}

One side effect is worth stating. A classifier that updates its history when
called is consulted only for transfers whose evaluation reaches one of its
leaves. In the default policy that leaf is in a rule of low priority, so the
order of rules decides which transfers the classifier sees.

\begin{proposition}[Linear cost]
\label{prop:linear}%
One decision costs time proportional to the total number of nodes in the
policy, plus one call to each classifier the policy names.
\end{proposition}

\begin{proof}
Each rule is evaluated at most once per decision, and within a rule each
node is visited at most once, since early stopping only skips nodes. A
comparison takes constant time, and a membership test takes time
proportional to the length of its list, which is part of the size of the
policy. Each classifier is called at most once per decision. Sorting is done
when the policy is loaded, not per decision.
\end{proof}

\section{A Controller and the Issue in Its Switching Rule}
\label{sec:controller}

\subsection{The Model}

The rule tree decides from thresholds on the current state and has no model
of how the two queues evolve. A third component, a controller derived from
optimal control, was designed to supply one. We model the two classes as two
queues that share one link, with a control that divides the link between
them:
\mathalt{Two equations. The rate of change of q H equals lambda H minus u times C times mu H.
  The rate of change of q C equals lambda C minus the quantity 1 minus u, times C times mu C.}{%
\begin{align}
  \dot q_H &= \lambda_H - u\, C \mu_H, \label{eq:qh}\\
  \dot q_C &= \lambda_C - (1-u)\, C \mu_C. \label{eq:qc}
\end{align}}
Here $q_H$ and $q_C$ are the numbers of operations that wait in the hot and
in the cold queue, and the dot denotes their rate of change in time. The
arrival rates of the two classes are $\lambda_H$ and $\lambda_C$, in
operations per second. $C$ is the capacity of the link, which the
implementation normalizes to $1$, and $\mu_H$ and $\mu_C$ are the service
rates of the two classes per unit of capacity, in operations per second. The
control $u$, a number between $0$ and $1$, is the share of the link given to
the hot queue. The cost to minimize is the weighted backlog over time, with
weight $\alpha$ on the hot queue and $\beta$ on the cold queue,
\looseness=-1 $\alpha > \beta$.

Pontryagin's principle \citep{pontryagin2018,liberzon2012} states that the optimal control minimizes a function, the Hamiltonian, at every instant. Here the
Hamiltonian is linear in $u$, and a linear function on an interval takes its
minimum at an end point. The optimal control is therefore always 0 or 1, a
``bang-bang'' control \citep{bellman1956}, and which one is decided by the
sign of a switching function,
\mathalt{S equals p H times C times mu H, minus p C times C times mu C}{%
\begin{equation}
  S = p_H\, C \mu_H - p_C\, C \mu_C ,
  \label{eq:switch}
\end{equation}}
in which $p_H$ and $p_C$ are the costs of one more operation of backlog in
the hot and in the cold queue, and $C$, $\mu_H$ and $\mu_C$ are as in
\cref{eq:qh,eq:qc}. When $S$ is positive the link should serve the hot
queue. The implementation approximates $p_H$ by $\alpha q_H$ and $p_C$ by
$\beta q_C$.

\subsection{The Inconsistency}

The code does not allocate link capacity but routes: for each new transfer it
computes $S$ and sends the transfer to the cold path when $S$ exceeds a small
dead band, which differs from the derivation in two ways.

The first is the sign. In the model, positive $S$ means that the hot queue
deserves the link, whereas the code responds to positive $S$ by sending new
work away from the hot path.

The second is the form. Choosing which queue a new transfer joins is a
different problem from choosing which queue to serve. A new transfer should
join the queue where it expects to wait less, and the expected wait is the
backlog divided by the service rate, so the rates should divide, not
multiply. A consistent rule sends a transfer to the cold path when
\mathalt{alpha times q H divided by mu H is greater than beta times q C divided by mu C}{%
\begin{equation}
  \frac{\alpha\, q_H}{\mu_H} > \frac{\beta\, q_C}{\mu_C},
  \label{eq:route}
\end{equation}}
that is, when the weighted expected wait in the hot queue, on the left,
exceeds that in the cold queue, on the right.

With the default parameters ($\alpha = 2$, $\beta = 1$, $\mu_H = 10{,}000$,
$\mu_C = 2{,}000$) the implemented rule sends a transfer cold whenever $q_H$
exceeds one tenth of $q_C$. \Cref{eq:route} would do so only when $q_H$
exceeds two and a half times $q_C$. The implemented rule therefore abandons
the hot path at 25 times less hot backlog than a consistent rule would.

The evidence agrees with this analysis. In the one recorded run on a device the controller made
1,000 decisions and chose cold every time (\cref{sec:res-loopback}). On a
synthetic grid of queue lengths the implemented rule chooses hot in 6 percent
of the cases and the consistent rule in 79 percent, and in a simulation with
Poisson arrivals the mean wait was 1.76 ms under the implemented rule and
0.28 ms under the consistent one. The simulation is synthetic, so its absolute numbers mean little, and only the
ordering of the two rules is the finding.

\subsection{Two Consistent Repairs}

There are two consistent repairs, which are to keep the controller as a router
and use \cref{eq:route}, or to make it allocate service between the queues
and correct the sign. The unit test of the controller pins the rule as
implemented (\cref{sec:tests}), and no result depends on it. The controller also has an extension that adds storage as a
third choice under GPU memory pressure by picking the smallest of three
costs, to which the same repair applies.

\fixnote{Decide before submission: apply one of the two repairs in the code and measure it (run R8), or keep this section as the analysis of the rule as implemented.}

\section{The Quantizer}
\label{ch:shrink}

\subsection{Scope of the Quantizer}

The method is not new. It follows TurboQuant \citep{zandieh2025turboquant}
in rotating a vector and rounding each coordinate, with a Hadamard matrix and
random signs in place of its random rotation matrix, and it takes the sign
sketch of QJL \citep{zandieh2024qjl}. The rounding is the classical scalar
quantizer of Lloyd and Max \citep{lloyd1982,max1960}, and the sketch rests on random projections
that preserve geometry \citep{johnson1984}. The contribution here is a packed wire format, an implementation whose sizes
are measured from the bytes written, and its placement, as
compression is applied to data in transit and at rest and never to the cache
on the GPU. On the serving path the codec is selected by configuration for a
whole run, and the compress field of the rule tree (\cref{sec:policy}) is the
means by which it can be selected per transfer.

Related work compresses the cache where it lives. KIVI quantizes it to 2 bits inside GPU memory and keeps a recent window at
full precision \citep{liu2024kivi}; transform coding compresses it for
compact storage \citep{staniszewski2026kvtc}; low-rank approximation reduces
its size in memory \citep{zhang2024lorc}; and a storage vendor has
reported using TurboQuant to speed up restoring caches from flash
\citep{alvarez2026flashblade}. No quality comparison against any of these was
run.

\subsection{The Four Steps}

\Cref{fig:quant-pipeline} shows the path of one cache vector $x$ of dimension
$d$ through the quantizer, and \cref{alg:quant} states the same four steps as
a procedure.

The first step normalizes. The length $\ell = \lVert x \rVert$ is stored as
a 32-bit number, and the work continues with the unit vector $u = x / \ell$.
The second step rotates. The unit vector is multiplied by a random diagonal
matrix $D$ of signs and then by a Hadamard matrix $H$, which gives
\mathalt{y equals 1 over the square root of d, times H times D times u}{$y = \tfrac{1}{\sqrt d} H D u$}.
The rotation preserves lengths and inner products, and it spreads the energy
of the vector so that every coordinate of $y$ has mean zero and variance
$1/d$. Cache vectors often have a few channels much larger than the rest, and
after the rotation they do not. The third step quantizes. Each coordinate of
$y$ is rounded to one of $2^b$ levels chosen to minimize the squared error
for that distribution, where $b$ is the number of bits per coordinate, so
that $b = 3$ gives eight levels. The fourth step sketches the residual $r$,
the difference between $x$ and its reconstruction $\hat x$. Its length
$\rho = \lVert r \rVert$ is stored together with the signs of $m$ random
projections of $r$, and from these signs the inner product of $r$ with any
query can be estimated, which corrects the attention score computed from the
reconstruction.

\begin{figure}[htbp]
  \centering
  \begin{tikzpicture}[diagram,
      alt={Flow diagram of the quantizer in six boxes joined by arrows. Top row, left to
        right. The input vector x enters box 1, normalize, which computes the length l of x
        and the unit vector u equals x divided by l, and stores l in 4 bytes. Box 2, rotate,
        computes y equals one over the square root of d, times H times D times u, and stores
        nothing. Box 3, quantize, finds for each coordinate y i the cell k i that contains it
        and replaces y i by the level c of k i; it stores the indices k in the ceiling of d
        times b over 8 bytes. Bottom row, right to left. Box reconstruct computes x hat equals
        l over the square root of d, times D times H times y hat, a step that the receiver
        repeats. Box residual computes r equals x minus x hat and its length rho, and stores rho
        in 4 bytes. Box 4, sketch, computes z j equals the sign of the inner product of s j
        with r divided by rho, for j from 1 to m, and stores the signs z in the ceiling of m
        over 8 bytes.},
      st/.style={tint=#1, text width=40mm, minimum height=19mm, font=\footnotesize, align=center, inner xsep=3pt, inner ysep=3pt}]
    \node[st=dBlue]   (n) at (0,0)     {1. Normalize\\$\ell = \lVert x \rVert,\quad u = x/\ell$\\store $\ell$: 4 bytes};
    \node[st=dBlue]   (r) at (5.1,0)   {2. Rotate\\$y = \tfrac{1}{\sqrt d}\, H D u$\\nothing stored};
    \node[st=dBlue]   (q) at (10.2,0)  {3. Quantize\\$k_i$: cell of $y_i$,\quad $\hat y_i = c_{k_i}$\\store $k$: $\lceil d b / 8 \rceil$ bytes};
    \node[st=dGray]   (h) at (10.2,-2.9) {Reconstruct\\$\hat x = \tfrac{\ell}{\sqrt d}\, D H \hat y$\\repeated by the receiver};
    \node[st=dOrange] (e) at (5.1,-2.9)  {Residual\\$r = x - \hat x,\quad \rho = \lVert r \rVert$\\store $\rho$: 4 bytes};
    \node[st=dOrange] (k) at (0,-2.9)    {4. Sketch\\$z_j = \operatorname{sign} \langle s_j, r/\rho \rangle$\\store $z$: $\lceil m / 8 \rceil$ bytes};
    \node[small text] (x) at (-3.0,0) {$x$};
    \draw[flow] (x) -- (n); \draw[flow] (n) -- (r); \draw[flow] (r) -- (q);
    \draw[flow] (q) -- (h); \draw[flow] (h) -- (e); \draw[flow] (e) -- (k);
  \end{tikzpicture}
  \caption[The four steps of the quantizer for one vector]{The four steps of the quantizer for one vector. $H$ is a Hadamard matrix, $D$ a random diagonal matrix of signs, $c$ the table of $2^b$ levels, and $s_1, \dots, s_m$ random sign vectors.}
  \label{fig:quant-pipeline}
\end{figure}

The proofs are in \cref{sec:proofs-quant}: the rotation is orthogonal; each rotated
coordinate has the stated mean and variance; the quantizer's levels satisfy
the conditions for a minimum; the sketch's estimate is unbiased for Gaussian
projections; and its standard deviation is at most
\mathalt{the square root of pi over 2 m, times the length of r times the length of q}{%
\begin{equation}
  \sqrt{\frac{\pi}{2m}}\; \lVert r \rVert\, \lVert q \rVert ,
  \label{eq:qjl}
\end{equation}}
where $m$ is the number of sign bits of the sketch, $\lVert r \rVert$ is the
length of the residual and $\lVert q \rVert$ is the length of the query. The
factor in front is 0.078 at $m = 256$. The implementation uses projections
with entries of plus and minus one instead of Gaussian entries, so the
estimate is approximately unbiased and the bound still holds.

\begin{pseudocode}{Packing and unpacking one cache vector}{alg:quant}
  \Function{Pack}{$x$} \Comment{$D$, the levels $c$ and $s_1, \dots, s_m$ are fixed by a seed}
    \State $\ell \gets \lVert x \rVert$ and $u \gets x / \ell$
    \State $y \gets \tfrac{1}{\sqrt d}\, H D u$ \Comment{fast transform, $d \log_2 d$ additions}
    \For{$i \gets 1$ to $d$}
      \State $k_i \gets$ index of the cell that contains $y_i$, and $\hat y_i \gets c_{k_i}$
    \EndFor
    \State $\hat x \gets \tfrac{\ell}{\sqrt d}\, D H \hat y$ \Comment{what the receiver will reconstruct}
    \State $r \gets x - \hat x$ and $\rho \gets \lVert r \rVert$
    \For{$j \gets 1$ to $m$}
      \State $z_j \gets \operatorname{sign} \langle s_j,\ r / \rho \rangle$
    \EndFor
    \State \Return $(\ell,\ \rho,\ k_1 \dots k_d,\ z_1 \dots z_m)$ \Comment{sizes in \cref{eq:packed}}
  \EndFunction
  \Function{Unpack}{$\ell$, $k_1 \dots k_d$}
    \State \Return $\tfrac{\ell}{\sqrt d}\, D H \hat y$ with $\hat y_i = c_{k_i}$
  \EndFunction
  \Function{Score}{$q$, $\hat x$, $\rho$, $z_1 \dots z_m$} \Comment{inner product with a query $q$}
    \State \Return $\langle q, \hat x \rangle + \rho\, \tfrac{\sqrt{\pi/2}}{m} \sum_{j=1}^{m} z_j \langle s_j, q \rangle$ \Comment{\cref{eq:qjl-est}}
  \EndFunction
\end{pseudocode}

\subsection{The Packed Format}

\Cref{fig:packed} shows one packed vector. Its size in bytes is
\mathalt{s packed equals 4 plus 4 plus the ceiling of d times b over 8, plus the ceiling of m
  over 8}{%
\begin{equation}
  s_{\text{packed}} = 4 + 4 + \left\lceil \frac{d\, b}{8} \right\rceil + \left\lceil \frac{m}{8} \right\rceil ,
  \label{eq:packed}
\end{equation}}
where the two terms of 4 bytes hold the length $\ell$ of the vector and the
length $\rho$ of its residual as 32-bit numbers, the third term holds the $d$
codes of $b$ bits each, the fourth holds the $m$ sign bits of the sketch, and
the brackets round up to a whole number of bytes.

\begin{figure}[htbp]
  \centering
  \begin{tikzpicture}[diagram,
      alt={Byte layout of one packed vector for 64 dimensions, 3 bits per coordinate and a
        256-bit sketch, drawn as a bar in four parts. Part one, 4 bytes: the length of the
        vector. Part two, 4 bytes: the length of the residual. Part three, 24 bytes: the 64
        three-bit codes. Part four, 32 bytes: the 256 sign bits of the sketch. A brace under
        the bar is labeled 64 bytes in total, against 128 bytes for the same vector in 16-bit
        floating point.},
      cell/.style={draw=dInk, line width=0.8pt, minimum height=11mm, font=\footnotesize, align=center, inner sep=2pt}]
    \node[cell, fill=dGray!12, minimum width=15mm] (a) at (0,0) {length\\4};
    \node[cell, fill=dGray!12, minimum width=15mm, anchor=west] (b) at (a.east) {residual\\4};
    \node[cell, fill=dBlue!12, minimum width=42mm, anchor=west] (c) at (b.east) {3-bit codes\\24};
    \node[cell, fill=dOrange!10, minimum width=54mm, anchor=west] (e) at (c.east) {sign sketch\\32};
    \draw[decorate, decoration={brace, amplitude=5pt, mirror}, draw=dGray, line width=0.6pt]
      ($(a.south west)+(0,-0.15)$) -- ($(e.south east)+(0,-0.15)$)
      node[midway, below=7pt, small text] {64 bytes (the same vector is 128 bytes in 16-bit floating point)};
  \end{tikzpicture}
  \caption[One packed vector for $d = 64$, $b = 3$, $m = 256$]{One packed vector for $d = 64$, $b = 3$, $m = 256$. Numbers are bytes.}
  \label{fig:packed}
\end{figure}

\Cref{tab:packed} gives sizes produced by the packing routine itself, for the
head dimension of the model used here and for a wider one.

\begin{table}[htbp]
  \centering
  \caption[Bytes per packed vector and the ratio to 16-bit floating point]{Bytes per packed vector, from the packing routine, and the ratio to 16-bit floating point.}
  \label{tab:packed}
  \small
  \begin{tabular}{L{2.6cm}L{3.0cm}L{3.0cm}L{2.4cm}L{2.4cm}}
    \toprule
    Dimension & Bits per coordinate & Sketch bits & Bytes & Ratio to 16-bit \\
    \midrule
    64 & 2 & 64 & 32 & 4.00 \\
    64 & 3 & 64 & 40 & 3.20 \\
    64 & 3 & 256 & 64 & 2.00 \\
    64 & 4 & 256 & 72 & 1.78 \\
    128 & 3 & 64 & 64 & 4.00 \\
    128 & 3 & 256 & 88 & 2.91 \\
    128 & 4 & 256 & 104 & 2.46 \\
    \bottomrule
  \end{tabular}
\end{table}

\subsection{Limitations}

Two limits of the implementation should be stated beside the result of
\cref{sec:res-quant}.

First, the sketch is not used where it travels. On the connector path the
receiver reconstructs each vector and writes it into the cache, and the
attention kernel then computes scores from the reconstruction alone. The 32
sketch bytes and the residual length cross the wire unused. Dropping them
would give 28 bytes per vector at $d = 64$, at the higher inner-product
error of \cref{tab:quant}, and applying the correction would need the
receiving kernel to use the sketch, which it does not.

Second, a sketch is only worth sending above a certain length. On synthetic
vectors with outlier channels the sketch reduced inner-product error only
from about 128 bits upward and made it worse below that.

\subsection{Analysis of the Quantizer}
\label{sec:proofs-quant}

Throughout, $d$ is the dimension of a cache vector and is a power of two,
which covers the head dimensions used here. The implementation pads other
dimensions with zeros to the next power of two. The inner product of two
vectors $u$ and $v$ is written $\langle u, v \rangle$ and is their dot
product (\cref{eq:dot}), the length of a vector is
$\lVert u \rVert = \sqrt{\langle u, u \rangle}$, and
$\operatorname{E}[\,\cdot\,]$ denotes the mean of a random quantity.

\subsubsection{The Rotation}

The Hadamard matrices are defined by $H_1 = (1)$ and
\mathalt{H of 2 n is the two by two block matrix with rows H n, H n and H n, minus H n}{%
\begin{equation}
  H_{2n} = \begin{pmatrix} H_n & H_n \\ H_n & -H_n \end{pmatrix},
\end{equation}}
so that $H_n$ is a square matrix with $n$ rows whose entries are all $+1$ or
$-1$. $D$ is a diagonal matrix with the diagonal entries
$\delta_1, \dots, \delta_d$, each of which is $+1$ or $-1$, and the rotation
is the matrix
\mathalt{R equals 1 over the square root of d, times H d times D}{%
\begin{equation}
  R = \frac{1}{\sqrt d}\, H_d\, D .
  \label{eq:rotation}
\end{equation}}

\begin{lemma}[The rotation is orthogonal]
\label[lemma]{lem:orth}%
Let $R^{\mathsf T}$ be the transpose of $R$ and $I$ the identity matrix.
Then $R R^{\mathsf T} = I$. Hence $\lVert R u \rVert = \lVert u \rVert$ and
$\langle R u, R v \rangle = \langle u, v \rangle$ for all vectors $u$ and
$v$, and $R^{\mathsf T}$ undoes $R$ exactly.
\end{lemma}

\begin{proof}
\emph{The Hadamard matrix.} First $H_d H_d^{\mathsf T} = d I$, by induction
on the size. The claim holds for $d = 1$. If $H_n H_n^{\mathsf T} = n I$,
multiplying the block form by its transpose gives the diagonal blocks
$H_n H_n^{\mathsf T} + H_n H_n^{\mathsf T} = 2 n I$ and the off-diagonal blocks
$H_n H_n^{\mathsf T} - H_n H_n^{\mathsf T} = 0$, which is the claim for $2n$.

\emph{The sign matrix.} $D D^{\mathsf T} = I$, because every diagonal entry
squares to $1$.

\emph{The product.} By \cref{eq:rotation} and the two previous steps,
\mathalt{R times R transposed equals 1 over d times H d times D times D transposed times H d
  transposed, which equals 1 over d times H d times H d transposed, which equals I}{%
\begin{equation}
  R R^{\mathsf T}
  = \frac{1}{d}\, H_d\, D D^{\mathsf T} H_d^{\mathsf T}
  = \frac{1}{d}\, H_d H_d^{\mathsf T} = I .
\end{equation}}

\emph{The consequences.} For a square matrix $R R^{\mathsf T} = I$ implies
$R^{\mathsf T} R = I$, so that
$\langle R u, R v \rangle = u^{\mathsf T} R^{\mathsf T} R\, v = \langle u, v \rangle$,
and the statement about lengths is the case $v = u$.
\end{proof}

\begin{lemma}[Every rotated coordinate has mean zero and variance $1/d$]
\label[lemma]{lem:moments}%
Fix a unit vector $u$, and let the diagonal entries
$\delta_1, \dots, \delta_d$ of $D$ be independent and equally likely to be
$+1$ or $-1$. Then each coordinate $y_i$ of $y = R u$ has mean $0$ and
variance $1/d$.
\end{lemma}

\begin{proof}
By \cref{eq:rotation}, coordinate $i$ of $y$ is
\mathalt{y i equals 1 over the square root of d, times the sum over j from 1 to d of H i j times
  delta j times u j}{%
\begin{equation}
  y_i = \frac{1}{\sqrt d} \sum_{j=1}^{d} H_{ij}\, \delta_j\, u_j ,
\end{equation}}
where $H_{ij}$ is the entry of $H_d$ in row $i$ and column $j$ and $u_j$ is
coordinate $j$ of $u$.

\emph{The mean.} Each $\delta_j$ has mean $0$, so $y_i$ has mean $0$.

\emph{The variance.} The $\delta_j$ are independent with $\delta_j^2 = 1$, so
the mean of $\delta_j \delta_k$ is $1$ if $j = k$ and $0$ otherwise, and
$H_{ij}^2 = 1$. Therefore
\mathalt{the mean of y i squared equals 1 over d times the sum over j and k of H i j times H i
  k times u j times u k times the mean of delta j delta k, which equals 1 over d times the
  sum over j of u j squared, which equals 1 over d}{%
\begin{equation}
  \operatorname{E}\big[y_i^2\big]
  = \frac{1}{d} \sum_{j,k} H_{ij} H_{ik}\, u_j u_k \operatorname{E}[\delta_j \delta_k]
  = \frac{1}{d} \sum_j u_j^2 = \frac{1}{d},
\end{equation}}
where the last step uses that $u$ is a unit vector, and because the mean is
$0$ this is the variance.
\end{proof}

The lemma gives two moments, not a distribution. Each $y_i$ is a sum of $d$
independent terms, and when no single coordinate of $u$ dominates, the
central limit theorem makes it close to a normal variable with that mean and
variance. That approximation is what justifies using one table of levels for
every coordinate of every vector.

\subsubsection{The Scalar Quantizer}

Let $Y$ be a random coordinate with density $f$ and finite variance. A
quantizer has $K = 2^b$ levels $c_1 < \dots < c_K$ and thresholds
$\theta_0 = -\infty < \theta_1 < \dots < \theta_{K-1} < \theta_K = \infty$,
and it replaces any value $y$ in cell $k$, the interval from $\theta_{k-1}$
to $\theta_k$, by the level $c_k$. Its mean squared error is
\mathalt{the mean squared error equals the sum over k from 1 to K of the integral from theta k
  minus 1 to theta k of the quantity y minus c k, squared, times f of y}{%
\begin{equation}
  \mathrm{MSE} = \sum_{k=1}^{K} \int_{\theta_{k-1}}^{\theta_k} (y - c_k)^2 f(y)\, dy ,
  \label{eq:mse}
\end{equation}}
in which term $k$ is the contribution of the values that fall into cell $k$.

\begin{proposition}[Conditions for a minimum \citep{lloyd1982,max1960}]
\label{prop:lloyd}%
If the levels and thresholds minimize the mean squared error of
\cref{eq:mse}, then each threshold is midway between its two levels,
$\theta_k = \tfrac12 (c_k + c_{k+1})$ for $k = 1, \dots, K-1$, and each level
is the mean of $Y$ over its cell,
$c_k = \operatorname{E}[\,Y \mid \theta_{k-1} < Y \le \theta_k\,]$ for
$k = 1, \dots, K$.
\end{proposition}

\begin{proof}
\emph{Thresholds, for fixed levels.} The error of a single value $y$ is
smallest when it is sent to the nearest level, and the values nearer to $c_k$
than to $c_{k+1}$ are those below the midpoint of the two. Any other
threshold sends some values to a level that is not their nearest and raises
the error.

\emph{Levels, for fixed thresholds.} The error is a sum of $K$ terms, each
of which depends on one level, and
\mathalt{the derivative of the mean squared error with respect to c k equals minus 2 times the
  integral from theta k minus 1 to theta k of the quantity y minus c k, times f of y}{%
\begin{equation}
  \frac{\partial\, \mathrm{MSE}}{\partial c_k}
  = -2 \int_{\theta_{k-1}}^{\theta_k} (y - c_k)\, f(y)\, dy ,
\end{equation}}
which is zero exactly when $c_k$ is the mean of $Y$ over the cell. The second
derivative is twice the probability of the cell, which is positive, so this
is the minimum for those thresholds.
\end{proof}

Lloyd's algorithm alternates the two conditions. Neither step can raise the
error, and the error is never negative, so it converges. The implementation runs
50 rounds on 200,000 samples of one coordinate of a random unit vector. For
$d = 64$ and $b = 3$ the resulting levels, multiplied by $\sqrt d$, are
$-2.118$, $-1.338$, $-0.764$, $-0.257$, $0.233$, $0.738$, $1.316$ and
$2.097$, computed from the code. On the sample they were fitted to, every
threshold is exactly the midpoint of its levels and every level is within
$0.0002$ of the mean of its cell, in units of the coordinate's standard
deviation. They are close to the classical eight levels for a normal
variable, $\pm 0.245$, $\pm 0.756$, $\pm 1.344$ and $\pm 2.152$
\citep{max1960}. The levels satisfy the conditions for the sample they were
fitted to; the conditions are necessary for a minimum and no more is claimed.

\subsubsection{The Sketch}

Let $r$ be the residual of one vector,
$\bar r = r / \lVert r \rVert$ the unit vector in its direction, and $q$ a
query. Let $s_1, \dots, s_m$ be random vectors of dimension $d$, and let the
sketch store the signs
$z_j = \operatorname{sign} \langle s_j, \bar r \rangle$, each of which is
$+1$ or $-1$. The estimate of the inner product $\langle q, r \rangle$ is
\mathalt{the estimate equals the length of r, times the square root of pi over 2, divided by m,
  times the sum over j from 1 to m of z j times the inner product of s j and q}{%
\begin{equation}
  \widehat{\langle q, r \rangle}
  = \lVert r \rVert\, \frac{\sqrt{\pi/2}}{m} \sum_{j=1}^{m} z_j\, \langle s_j, q \rangle ,
  \label{eq:qjl-est}
\end{equation}}
in which the hat marks the estimate, $\lVert r \rVert$ is the stored length
of the residual, and $\langle s_j, q \rangle$ is the projection of the query
on the same random vector that produced the sign $z_j$.

\begin{lemma}[The mean of a sign times a normal variable]
\label[lemma]{lem:sign}%
Let $X$ and $Y$ be jointly normal random variables with mean $0$, let
$\sigma^2 > 0$ be the variance of $X$, and let $\kappa$ be the covariance of
$X$ and $Y$. Then
$\operatorname{E}[\operatorname{sign}(X)\, Y] = \sqrt{2/\pi}\; \kappa / \sigma$,
where $\operatorname{sign}(X)$ is $+1$ if $X$ is positive and $-1$
otherwise.
\end{lemma}

\begin{proof}
\emph{Splitting $Y$.} Write $Y = a X + V$ with $a = \kappa / \sigma^2$, which
defines the variable $V = Y - a X$. Then $V$ has mean $0$ and is uncorrelated
with $X$, since their covariance is $\kappa - a \sigma^2 = 0$, and for
jointly normal variables uncorrelated means independent. Hence
$\operatorname{E}[\operatorname{sign}(X) V] = \operatorname{E}[\operatorname{sign}(X)] \operatorname{E}[V] = 0$
and
$\operatorname{E}[\operatorname{sign}(X) Y] = a \operatorname{E}\lvert X \rvert$.

\emph{The mean of $\lvert X \rvert$.} For a normal variable with mean $0$ and
variance $\sigma^2$,
\mathalt{the mean of the absolute value of X equals 2 over sigma times the square root of 2 pi,
  times the integral from zero to infinity of x times the exponential of minus x squared
  over 2 sigma squared, which equals sigma times the square root of 2 over pi}{%
\begin{equation}
  \operatorname{E}\lvert X \rvert
  = \frac{2}{\sigma \sqrt{2\pi}} \int_0^\infty x\, e^{-x^2 / 2\sigma^2}\, dx
  = \sigma \sqrt{\frac{2}{\pi}} ,
\end{equation}}
and $a \sigma \sqrt{2/\pi} = \sqrt{2/\pi}\; \kappa / \sigma$.
\end{proof}

\begin{proposition}[The estimate is unbiased for normal projections \citep{zandieh2024qjl}]
\label{prop:unbiased}%
If the vectors $s_1, \dots, s_m$ are independent and their entries are
independent standard normal variables, then the mean of the estimate of
\cref{eq:qjl-est} is $\langle q, r \rangle$.
\end{proposition}

\begin{proof}
Fix $q$ and $r$, and write $X_j = \langle s_j, \bar r \rangle$ and
$Y_j = \langle s_j, q \rangle$ for the two projections on $s_j$, so that
$z_j = \operatorname{sign}(X_j)$. The pair $X_j$, $Y_j$ is jointly normal
with mean $0$, $X_j$ has variance $\lVert \bar r \rVert^2 = 1$, and the
covariance of the pair is $\langle \bar r, q \rangle$. By \cref{lem:sign}
with $\sigma = 1$ and $\kappa = \langle \bar r, q \rangle$,
$\operatorname{E}[z_j Y_j] = \sqrt{2/\pi}\, \langle \bar r, q \rangle$. Each
of the $m$ terms of the sum has this mean, so the mean of the estimate is
\mathalt{the length of r, times the square root of pi over 2, times the square root of 2 over
  pi, times the inner product of r bar and q, which equals the inner product of q and r}{%
\begin{equation}
  \lVert r \rVert\, \sqrt{\frac{\pi}{2}}\, \sqrt{\frac{2}{\pi}}\, \langle \bar r, q \rangle
  = \lVert r \rVert\, \langle \bar r, q \rangle = \langle q, r \rangle .
\end{equation}}
\end{proof}

\begin{proposition}[Standard deviation of the estimate]
\label{prop:qjlvar}%
If the vectors $s_1, \dots, s_m$ are independent and the entries of each are
uncorrelated with mean $0$ and variance $1$, the standard deviation of the
estimate of \cref{eq:qjl-est} is at most the bound of \cref{eq:qjl}.
\end{proposition}

\begin{proof}
\emph{One term.} With $Y_j = \langle s_j, q \rangle$ as before, the terms
$z_j Y_j$ are independent and identically distributed. Since $z_j^2 = 1$,
the variance of one term is at most
$\operatorname{E}[Y_j^2] = \lVert q \rVert^2$, where the last step uses only
that the entries of $s_j$ are uncorrelated with mean $0$ and variance $1$.

\emph{The estimate.} The variance of the mean of $m$ such terms is at most
$\lVert q \rVert^2 / m$, and the estimate is that mean multiplied by
$\lVert r \rVert \sqrt{\pi / 2}$. Its variance is therefore at most
$\tfrac{\pi}{2m} \lVert r \rVert^2 \lVert q \rVert^2$, and the square root of
this is the bound of \cref{eq:qjl}. For $m = 256$ the factor
$\sqrt{\pi / 512}$ is $0.078$.
\end{proof}

The implementation draws the entries of each $s_j$ as $+1$ or $-1$ with
equal probability. Those entries are independent with mean $0$ and variance
$1$, so \cref{prop:qjlvar} applies as stated. \Cref{prop:unbiased} does not:
$X_j$ and $Y_j$ are then sums of $d$ independent terms and only approximately
jointly normal, so the estimate is approximately unbiased, with an error
that shrinks as $d$ grows. The size of that bias at $d = 64$ was not
measured.

\section{The Storage Tier}
\label{sec:tier}

\subsection{Place in the System}

Serving systems are converging on a hierarchy for cache: GPU memory, host
memory, local flash, and a shared file system below that
\citep{hsu2026fourtier,hsu2026kv,ibm2025storage}. We implement one step of it, from memory to a local NVMe drive, as a place to
which the rule tree can send a cold block when the GPU is nearly full. The
tier is exercised by its own tests and by the sweep of \cref{sec:res-nvme},
and no experiment in this work spills a block of a running request.

\subsection{Design}

The tier keeps one backing file. A block is stored under its key in a run of
4 KiB extents found by a first-fit search of a free list; freed runs are
merged with their neighbors. The tier keeps a 32-bit checksum of each block
in its index and verifies it on every read, so a damaged block is reported
and never returned as data. The index is held in memory, so the tier is
spill space for one run of the server and not a store that survives a
restart. Writes are batched: the tier queues blocks and writes a
batch with one call followed by one flush to stable storage. If the file
system allows it, the file is opened for direct I/O, which bypasses the
operating system's cache and requires every length, offset and buffer address
to be a multiple of 4,096. Where direct I/O is refused, the tier falls back
to buffered writes and records why. The procedures are stated with the
analysis that follows, which proves the alignment conditions, the invariants
\looseness=-1 of the free list, and the classes of error that the checksum detects.

\subsection{Analysis of the Storage Tier}
\label{sec:proofs-storage}

\Cref{alg:tier,alg:tier-read} give the write and the read path of the tier,
and \cref{alg:alloc}, further below, the allocator beneath them. The tier
works in units of
$A = 4{,}096$ bytes, the alignment that direct I/O requires. A block of
$\var{len}$ bytes occupies a slot of $w = \lceil \var{len} / A \rceil\, A$
bytes, that is, its length rounded up to a multiple of $A$. The tier keeps
the end of the file, written $\var{end}$, a list of free extents, each with
an offset and a size, and an index from each key to the offset, length, slot
\looseness=-1 size and checksum of its block.

\begin{pseudocode}{Spilling a block to the storage tier}{alg:tier}
  \Function{Spill}{$\var{key}$, $\var{block}$}
    \State append $(\var{key}, \var{block})$ to the pending list
    \If{the pending list holds $n$ blocks} \Comment{$n$ is the batch size}
      \State \Call{WriteBatch}{the pending list} on a worker thread, and empty the list
    \EndIf
  \EndFunction
  \Function{WriteBatch}{blocks $1, \dots, n$ of $\var{len}_1, \dots, \var{len}_n$ bytes}
    \State $w_i \gets \lceil \var{len}_i / A \rceil\, A$ for each $i$ \Comment{slot of block $i$, $A = 4{,}096$}
    \State $\var{base} \gets \Call{Allocate}{w_1 + \dots + w_n}$ \Comment{\cref{alg:alloc}}
    \State place each block at the start of its slot in one aligned buffer, padded with zeros
    \State write the buffer at offset $\var{base}$ with one call, then flush to stable storage once
    \For{$i \gets 1$ to $n$}
      \State $\var{index}[\var{key}_i] \gets$ (offset of slot $i$,\ $\var{len}_i$,\ $w_i$,\ \Call{Crc32}{block $i$})
    \EndFor
  \EndFunction
\end{pseudocode}

\begin{pseudocode}{Reading a block back from the storage tier}{alg:tier-read}
  \Function{Retrieve}{$\var{key}$}
    \State \textbf{if} the block is still pending \textbf{then return} it from memory
    \State $(o,\ \var{len},\ w,\ \var{crc}) \gets \var{index}[\var{key}]$
    \State read the slot of $w$ bytes at offset $o$ and keep its first $\var{len}$ bytes
    \State \textbf{if} \Call{Crc32}{bytes kept} differs from $\var{crc}$ \textbf{then fail} with an integrity error
    \State \Return the bytes kept
  \EndFunction
\end{pseudocode}

\begin{proposition}[Every transfer meets the conditions of direct I/O]
\label{prop:align}%
In every write and read that the tier issues with direct I/O, the file
\looseness=-1 offset, the length and the address of the buffer are multiples of $A$.
\end{proposition}

\begin{proof}
\emph{Lengths.} A block of $\var{len}$ bytes is given a slot of
$w = \lceil \var{len} / A \rceil\, A$ bytes. A batch is written as one run of
such slots, whose length is a sum of multiples of $A$, and a read covers one
\looseness=-1 slot or a run of them.

\emph{Offsets.} The argument is an induction over the operations of the
allocator (\cref{alg:alloc}). The end of the file starts at $0$. An allocation returns either
the end of the file, which then advances by a multiple of $A$, or the start
of a free extent, which is then shortened from the front by a multiple of
$A$. A release returns a range that an allocation handed out, split at slot
boundaries, so its offset and size are multiples of $A$, and merging two such
ranges gives another. Every offset and size that the allocator ever holds is
therefore a multiple of $A$.

\emph{Addresses.} Buffers for direct I/O are obtained from anonymous memory
maps, which begin on a page boundary, and the page size on these machines is
$4{,}096$ bytes.
\end{proof}

The proposition assumes that $A$ is a multiple of the logical block size of
the device, which is 512 or 4,096 bytes on the drives used here.

\begin{pseudocode}{Allocating and releasing extents}{alg:alloc}
  \Function{Allocate}{$w$} \Comment{$w$ bytes, a multiple of $A$}
    \For{each free extent of $g$ bytes at offset $o$, in order of offset} \Comment{first fit}
      \If{$g \ge w$}
        \State shorten the extent to $g - w$ bytes at offset $o + w$, or remove it if $g = w$
        \State \Return $o$
      \EndIf
    \EndFor
    \State $o \gets \var{end}$, then $\var{end} \gets \var{end} + w$ and \Return $o$ \Comment{grow the file}
  \EndFunction
  \Function{Release}{$o$, $w$} \Comment{$w$ bytes at offset $o$}
    \State insert the extent of $w$ bytes at offset $o$ into the free list and sort it by offset
    \State merge every extent with a neighbor that begins where it ends
    \State \textbf{if} the last extent ends at $\var{end}$ \textbf{then} set $\var{end}$ to its offset and remove it
  \EndFunction
\end{pseudocode}

\begin{proposition}[No two blocks share space]
\label{prop:freelist}%
At all times the free extents are sorted and pairwise disjoint, no two of
them are adjacent, none overlaps an allocated range, and all of them end
below the end of the file. Consequently two allocated ranges never overlap.
\end{proposition}

\begin{proof}
The properties named in the first sentence form an invariant, which is
shown to hold at the start and after every operation of \cref{alg:alloc}.

\emph{At the start.} There is no free extent and no allocated range, so all
properties hold.

\emph{An allocation.} The range is taken from the front of a free extent, or
from the old end of the file upward. In the first case the range was free, so
by the invariant it overlapped no allocated range and no other free extent,
and the part of the extent that remains keeps its end, so the list stays
sorted, disjoint and without adjacent pairs. In the second case the range
lies above everything, and the end of the file moves up past it.

\emph{A release.} A release adds back a range that was allocated, which by
the invariant overlapped no free extent, so the free extents stay disjoint.
The release then sorts the list, merges every pair in which one extent ends
where the next begins, and, if the last extent ends at the end of the file,
moves the end of the file down to the start of that extent and drops it.
Merging and dropping do not change which bytes below the new end are free.

\emph{Conclusion.} In both cases of an allocation the new range is disjoint
from every allocated range, which gives the final statement.
\end{proof}

Merging is what allows reuse: two adjacent freed slots of 4 KiB become one
extent of 8 KiB that can hold a block neither could hold alone.

\begin{proposition}[What the checksum detects]
\label{prop:crc}%
Let a block be read back with some of its bits changed. The 32-bit cyclic
redundancy check used by the tier reports the change in each of these cases:
one bit is changed; two bits are changed, in a block shorter than
$2^{32} - 1$ bits; or all changed bits lie within four consecutive bytes. If
instead every pattern of changed bits over the block is equally likely, the
probability of a change that goes unreported is less than $2^{-32}$.
\end{proposition}

\begin{proof}
\emph{Setting.} Treat the block as a polynomial in a variable $x$ over the
field with two elements, with one coefficient per bit, so that the bit at
position $i$ is the coefficient of $x^i$. The check is the remainder of the
block, multiplied by $x^{32}$, on division by a fixed generator polynomial
$G$ of degree 32; the constants that the standard adds at the start and the
end are the same for two blocks of equal length and cancel when they are
compared. Changing a set of bits adds an error polynomial $E$ to the block,
whose nonzero coefficients mark the changed bits. The change goes unreported
exactly when the remainder is unchanged, that is, when $G$ divides
$x^{32} E$. The constant term of $G$ is $1$, so $G$ and $x$ have no common
factor, and this happens exactly when $G$ divides $E$.

\emph{One changed bit.} The error is $E = x^i$, which $G$ cannot divide,
because $G$ has no common factor with $x$ and has degree 32.

\emph{Changes within four consecutive bytes.} The changed bits lie within 32
consecutive positions, whatever the order of bits inside a byte. They give
$E = x^i E'$, where $E'$ is nonzero and of degree at most 31, and $G$ divides
neither factor.

\emph{Two changed bits.} The error is $E = x^i (x^k + 1)$, where $k \ge 1$
is the distance between the two bits, and $G$ divides it only if $G$ divides
$x^k + 1$. The smallest such $k$ is the order of $x$ modulo $G$. For this
generator the order is $2^{32} - 1$: a direct computation shows that
$x^{2^{32}-1}$ leaves remainder $1$ and that $x^{(2^{32}-1)/p}$ does not for
each prime $p$ that divides $2^{32} - 1$, namely 3, 5, 17, 257 and 65,537. In
a block shorter than $2^{32} - 1$ bits, about 512 MiB, two changed bits are
less than that far apart.

\emph{Random changes.} In a block of at least 32 bits, sending a pattern of
changed bits to its remainder on division by $G$ is a linear map onto the
$2^{32}$ possible remainders, so every remainder is the image of equally many
patterns. The unreported patterns are those with remainder zero, a fraction
$2^{-32}$ of all patterns, and one of them is the pattern with no change at
all.
\end{proof}

The generator has 15 terms. A polynomial that is a multiple of $x + 1$ has an
even number of terms, so $x + 1$ is not a factor of this generator, and the
check does not detect every change of an odd number of bits. The checksum is kept in the tier's index in memory,
not beside the block on the drive, so it protects against a block that comes
back different from what was written, and not against loss of the index.

\subsection{Batching and a Model of Its Effect}

A flush has a cost that does not depend on how many bytes were written before
it, and writing several blocks per flush spreads that cost. The simplest
\looseness=-1 model is a fixed time per batch and a time per block:
\mathalt{t batch of n equals t fix plus n times t block, and X of n equals n times s block
  divided by the quantity t fix plus n times t block, which approaches s block over t block
  as n grows}{%
\begin{equation}
  t_{\text{batch}}(n) = t_{\text{fix}} + n\, t_{\text{blk}},
  \qquad
  X(n) = \frac{n\, s_{\text{blk}}}{t_{\text{fix}} + n\, t_{\text{blk}}}
  \;\longrightarrow\; \frac{s_{\text{blk}}}{t_{\text{blk}}}
  \quad\text{as } n \text{ grows},
  \label{eq:coalesce}
\end{equation}}
where $n$ is the number of blocks in a batch, $t_{\text{batch}}(n)$ is the
time to write and flush the batch, $t_{\text{fix}}$ is the part of that time
which does not depend on $n$ and includes the flush, $t_{\text{blk}}$ is the
time added by each block, $s_{\text{blk}}$ is the size of a block in bytes,
and $X(n)$ is the resulting write throughput in bytes per second.
\Cref{sec:res-nvme} fits this to the measured sweep and finds it adequate for
batches of ten or more and poor for single writes. The batch size is a tuning value, since a larger batch also delays
the moment a block is safely on disk, and the read results in that section
suggest very large batches have a cost of their own.

This completes the description of the components. The next chapter
describes the machines on which they run and the method by which they are
measured.
